\chapter{Modell und Training}
\label{ch:training}

% ==============================================================================
\section{Ablauf vom Datensatz zum Modell}
\label{sec:train-ablauf}

\Cref{fig:kette} zeigt die Kette von der Aufnahme bis zur Ausführung. Jeder Schritt ist ein
eigenes Programm, das auch aus der Bedienoberfläche aufgerufen wird, und jeder Schritt
schreibt seine Einstellungen in die Ausgabe. Der exportierte Datensatz ist eigenständig: Er
kann auf einen anderen Rechner kopiert und dort trainiert werden, ohne dass dort die
Hardware-Bibliotheken installiert sind.

\begin{figure}[htbp]
  \centering
  \begin{tikzpicture}[font=\small, node distance=4mm and 6mm,
      st/.style={draw, rounded corners=1.5mm, fill=blue!8, align=center, minimum height=12mm,
                 text width=23mm},
      da/.style={draw, fill=black!5, align=center, minimum height=9mm, text width=23mm,
                 font=\scriptsize},
      arr/.style={-{Stealth[length=2mm]}, thick}]
    \node[st] (rec) {Aufnahme\\\code{record}};
    \node[st, right=of rec] (exp) {Export\\\code{export}};
    \node[st, right=of exp] (tr) {Training\\\code{train}};
    \node[st, right=of tr] (inf) {Ausführung\\\code{infer}};
    \draw[arr] (rec) -- (exp); \draw[arr] (exp) -- (tr); \draw[arr] (tr) -- (inf);
    \node[da, below=of rec] {Episoden\\npz + json};
    \node[da, below=of exp] {LeRobot-Datensatz\\+ Herkunft};
    \node[da, below=of tr] {Policy\\+ \code{bc_policy.json}};
    \node[da, below=of inf] {Fahrtprotokolle\\+ Auswertung};
    \node[draw, densely dashed, rounded corners=1.5mm, fit=(rec)(inf), inner sep=3mm,
          label={[font=\scriptsize]above:Prüfung der Einstellungen bei jedem Übergang}] {};
  \end{tikzpicture}
  \caption[Kette von der Aufnahme bis zur Ausführung]{Kette von der Aufnahme bis zur
  Ausführung. Bei jedem Übergang wird geprüft, ob Schema, Takt, Kameras, Tempo und Override
  zusammenpassen.}
  \label{fig:kette}
\end{figure}

\paragraph{Prüfung beim Export.} Der Export bricht ab, wenn Schemaversion, Takt oder Kameras
zwischen den Quellen abweichen oder wenn Override, Sollwertrate oder eine Geschwindigkeit des
Planers nicht in allen Episoden gleich sind. Das Tempo wird mitgelernt, und ein Datensatz aus
gemischten Geschwindigkeiten würde eine Policy mit unbestimmtem Tempo ergeben. Verworfene und
als misslungen markierte Episoden gehen nie ein.

% ==============================================================================
\section{Netzarchitektur}
\label{sec:netz}

Verwendet wird die faltungsbasierte Diffusion Policy aus LeRobot (\cref{sec:dp},
\cref{fig:dp-netz}). \Cref{tab:hyper} nennt die Einstellungen, die folgenden Abschnitte
begründen die wichtigsten.

\begin{table}[htbp]
  \centering
  \caption{Einstellungen der Diffusion Policy}
  \label{tab:hyper}
  \small
  \begin{tabularx}{\textwidth}{@{}l X@{}}
    \toprule
    Parameter & Wert \\
    \midrule
    Bildkodierer & ResNet18 je Kamera, ImageNet-Gewichte, Spatial Softmax \\
    U-Net-Kanäle & 256, 512, 1024 (\num{89.5} Millionen Parameter gesamt) \\
    Beobachtungen $T_o$ & 2 Takte \\
    Vorhersagehorizont $T_p$ & 16 Aktionen, davon $T_a = 8$ nutzbar \\
    Neuvorhersage & alle 2 Takte, überlappende Folgen ungewichtet gemittelt \\
    Entrauschen & 100 Schritte im Training, 10 \acs{ddim}-Schritte bei der Ausführung \\
    Vorhersageziel & Aktionsfolge direkt (\cref{eq:loss-x0}) \\
    Bildaugmentierung & zufälliger Ausschnitt \SI{90}{\percent} im Training, Mitte bei der
      Ausführung \\
    Optimierer & AdamW~\cite{loshchilov2019}, Lernrate $10^{-4}$, Cosinus-Verlauf \\
    Gewichte & gleitender Mittelwert (\acs{ema}, \num{0.999}), gespeichert werden die
      gemittelten Gewichte \\
    Genauigkeit & bf16 \\
    \bottomrule
  \end{tabularx}
\end{table}

\paragraph{Größe des U-Nets.} Die Kanalzahlen sind kleiner als die Vorgabe von LeRobot. So
passt das Training in die \SI{16}{\giga\byte} der Grafikkarte, und die Vorhersage bleibt auf
einem Laptop ausführbar. Auf der RTX~5070~Ti dauert ein Schritt mit Batchgröße 64 etwa
\SI{0.21}{\second} bei \SI{5.6}{\giga\byte} Speicher; das Laden der Daten ist mit
\SI{0.02}{\second} je Schritt nicht der Engpass.

\paragraph{Horizont.} Zwei Beobachtungen geben der Policy eine Information über die Richtung
der Bewegung, die ein einzelner Zustand nicht enthält. 16 vorhergesagte Aktionen entsprechen
gut einer Sekunde und glätten die Bewegung, 8 davon werden höchstens ausgeführt, bevor eine
neuere Vorhersage vorliegt.

\paragraph{Vorhersageziel.} Zunächst wurde die Vorgabe von LeRobot verwendet: Das Netz sagt
das enthaltene Rauschen voraus (\cref{eq:loss-eps}).
\begin{description}
  \item[Problem] Mit nur zehn Entrauschungsschritten waren die erzeugten Aktionsfolgen
    gezackt. Der größte Sprung innerhalb einer Folge lag bei \SIrange{0.10}{0.12}{rad},
    in den Trainingsdaten bei \SI{0.015}{rad}. Die Sicherheitsprüfung stoppte die Fahrt schon
    in den ersten Takten. Mit 100 Schritten war die Folge glatt, die Vorhersage dauerte dann
    aber etwa \SI{600}{\milli\second}.
  \item[Lösung] Das Netz sagt die Aktionsfolge direkt vorher (\cref{eq:loss-x0}). Damit
    entsprechen die Sprünge denen der Trainingsdaten. Die Validierung protokolliert seitdem
    den Sprung innerhalb der Folge neben dem Wert der Daten.
\end{description}

\paragraph{Augmentierung.} Der zufällige Bildausschnitt deckt kleine Verschiebungen der
Draufsichtkamera ab. Beim Fischauge der Handgelenkkamera ist sein Nutzen geringer, weil ein
Ausschnitt das optische Zentrum verschiebt, was keiner echten Kamerabewegung entspricht.

% ==============================================================================
\section{Training und Validierung}
\label{sec:validierung}

Ein Teil der Episoden wird für die Validierung zurückgehalten, und zwar ganze Episoden, nicht
einzelne Takte, weil benachbarte Takte fast gleich sind und die Validierung sonst zu
optimistisch wäre. In festen Abständen protokolliert das Training neben dem Diffusionsverlust
drei anschaulichere Größen: den mittleren Gelenkfehler der erzeugten Aktionsfolgen in rad,
den Anteil richtig vorhergesagter Greiferzustände und den größten Sprung innerhalb einer
Folge.

Mit dem Modell wird eine Beschreibungsdatei gespeichert (\code{bc_policy.json}). Sie enthält
Schema, Takt, Kameras, LeRobot-Version, Override, Sollwertrate und Tempo der Aufnahme, Start-
und Endstellung, eine Referenzbahn, den Tisch und die Hülle der Demonstrationen. Die
Ausführung liest daraus ihre Prüfungen ab (\cref{ch:ausfuehrung,ch:sicherheit}).

\paragraph{Problem: eine Uhr im Bild.} Die erste simulierte Kamera schrieb zur Kontrolle einen
Bildzähler ins Bild. Die Policy las ihn als Zeit: Die erste Fahrt gelang, danach fuhr sie direkt
ans Bahnende. Platzhalterbilder sind seitdem statisch. Der Fall zeigt allgemein, dass die
Policy jede Information nutzt, die mit dem Fortschritt der Aufgabe zusammenhängt, auch
unbeabsichtigte.

% ==============================================================================
\section{Prüfung eines Modells ohne Anlage}
\label{sec:replay}

Ob ein Modell brauchbar ist, lässt sich in zwei Stufen prüfen, bevor es an die Anlage geht.

\begin{description}
  \item[Stufe 1, offline] Bilder und Zustand einer aufgezeichneten Episode gehen Takt für Takt
    in die Policy, die Vorhersage wird mit der aufgezeichneten Aktion verglichen. Ist der
    Fehler schon hier groß, liegt das Problem bei Daten oder Training.
  \item[Stufe 2, \acs{vm}] Die Policy fährt die \ac{vm}. Die Bilder kommen aus der Aufnahme, je
    Takt wird ein Bild weitergeschaltet, der Zustand kommt vom Roboter. Verglichen wird mit der
    idealen Bahn der Aufnahme.
\end{description}
Die Grenze von Stufe~2 ist, dass die Bilder nicht auf den Arm reagieren. Weicht er ab, sieht die
Policy trotzdem die Bilder der Aufnahme. Die Stufe eignet sich deshalb zum Vergleich von
Modellen, nicht für absolute Aussagen über die Anlage. Auf diesem Weg wurden mehrere
Varianten von Daten und Training verglichen (\cref{ch:versuche}).
